\section{Complete Form Scores and Operational Comparisons}\label{app:operational}
\subsection{Shared-Coordinate Corrections and the Complete Score Space}
\label{app:fullscores}
Fix a menu of coordinate sets $\mathcal T$. For binary features write
$g_{T,y}(x_T)=\sum_{S\subseteq T}\gamma_{T,y,S}\chi_S(x)$.
Uniqueness of the Fourier expansion gives the following elementary
representation.

\begin{proposition}[All shared-coordinate corrections]\label{prop:sharing}
The score tables satisfy $\sum_Tg_{T,y}(x_T)=yd(x)$ on the full binary cube
if and only if
\[
       \sum_{T\supseteq S}\gamma_{T,y,S}=y\widehat d(S)
       \quad\text{for every }S.
\]
For a fixed positive design $q$, every change to a globally unbiased score
is a sum of pairwise transfers of functions on shared coordinates:
add $f(X_{T\cap T'},Y)/q_T$ to form $T$ and subtract the same function
divided by $q_{T'}$ from form $T'$.
\end{proposition}
\begin{proof}
Equate the coefficients of each $\chi_S$. For the second claim, any score
difference has, for each $S$ and each label, coefficients summing to zero
over forms containing $S$. Choose one of those forms as an anchor and
transfer every other coefficient to or from the anchor. Each transfer is
observable under both forms. Terms appearing in only one form must have
zero difference. Summing the transfers proves the statement.
\end{proof}
Thus including all such terms uses the entire unbiased score space on the
supplied menu. This is a product-basis implementation of standard control
variates and known-coarsening influence-function optimization, not a new
general missing-data efficiency bound.

The implementation works in cell values. A block of the constraint matrix
is the normalized Walsh transform of a form's table, placed in the rows of
its contained coordinate subsets. If $F$ is the resulting matrix, the
constraints are $Fg_y=y\widehat d$. Its rows are orthogonal:
$FF^{\mathsf T}$ is diagonal, with entry
$\sum_{T\supseteq S}2^{-|T|}$ for subset $S$. We use this identity to
repair numerical equality residuals explicitly. Rotated second-order cones
represent $g_{T,y}(u)^2\le q_Tv_{T,y}(u)$.

For a given population, marginal form--label frequencies are sufficient for
the objective. No full-profile enumeration is needed. With $L$ forms of
size at most $b$, there are at most $2L2^b$ score entries and $2L2^b$
coefficient constraints. A tree provides the target coefficients through
its leaves, or by Fourier transforming its known completion-score
decomposition. Choosing an unrestricted menu can still require
$\binom mb$ forms; this implementation does not solve that combinatorial
design problem.

For ordinal features we split each coordinate at all thresholds used for it
in the frozen tree, with a fixed middle cut for otherwise unused padding
coordinates. We use orthonormal Helmert contrasts on those categories and
their tensor products. The same argument is exact on the resulting product
space. This retains every observed tree-threshold category, but not
within-category distinctions between raw ordinal answers. Its cell count
is $\sum_T\prod_{j\in T}s_j$, where $s_j$ is the category count.

\paragraph{Calibration and certificates.}
The larger-tree experiment uses a model fitting sample of
800, the next 500 complete calibration records, and the disjoint remaining
evaluation sample. A common 350/150 internal calibration split chooses
among uniform joint-law shrinkage weights $0.02,0.1,0.3,1$ and the prefix
baseline. Selection minimizes validation second moment, averaged over
random form assignment. The selected method is refitted on all 500 records.
The fitted design uses $q_T\ge0.1/L$. Evaluation records choose no
hyperparameter or score.

The menu oracle separately optimizes over every score unbiased on the full
product profile space and all design probabilities on that same menu,
without the probability floor.
It is a benchmark, not a fitted method. We reconstruct feasible primal
tables, and a dual witness provides its lower bound. For SUSHI this
full-product unbiasedness requirement is more conservative than restricting
it to ranking-feasible profiles; the eight-feature comparison separately
examines both supports. For a quadratic
objective with cell masses $p$, a witness $h_y$ on the coefficient constraints
has information
$I_T=\sum_{y,u}(F_T^{\mathsf T}h_y)_u^2/p_{T,y,u}$.
With unrestricted probabilities the lower bound on the second moment is
$(\sum_y h_y^{\mathsf T}y\widehat d)^2/\max_TI_T$.
If $p_{T,y,u}=0$, the corresponding numerator must vanish; the code projects
the witness onto those linear zero-cell constraints. Numerical bounds and
equality residuals are retained for every task. This oracle quantifies
the restriction to the chosen menu; it does not lower-bound the optimum
over all possible fixed forms.

The largest individual variance-ratio bracket width is 0.00177 on the
binary tasks and 0.01639 on ordinal tasks. The displayed oracle summaries
use the upper endpoints; both endpoints are released. These are numerical
floating-point certificates, not interval-arithmetic proofs of exact
real-number optima.

\subsection{Independent Model Versions on Eight Features}\label{app:medium}
For each of the 45 SUSHI and 25 binary SAPA primary-split tasks, we select
eight target-excluding features at equally spaced indices of the existing
feature list. A depth-two old classifier is fitted on the first 400 records;
a depth-three new classifier is independently fitted on the next 400.
Each uses minimum leaf size 20 and a fixed seed. The remaining records
define the evaluation law. We examine both probability outputs and
deterministic threshold outputs.

The audit first follows the old tree, then the new tree, reusing an answer
when both trees ask the same feature. Its cap is the maximum size of the
union of the two reached paths over all 256 profiles, at most five. This
cap is computed from the frozen rules, not chosen using the evaluation
population. Hence the branching coefficient is $\|d\|_\infty^2$ and its
population variance is $\Var_P(Yd)$. We optimize all fixed forms at that
same cap, including every possible answer table. Both coefficients and
population-specific variances have independently evaluated primal/dual
bounds.

For SUSHI, we additionally enumerate all $8!$ orders of the eight nontarget
alternatives and retain only feasible profiles of the eight comparisons.
The optimization is repeated on this ranking support, while retaining the
same audit cap. This checks the effect of logical restrictions among
pairwise comparisons. The target pair involves the two other alternatives,
so either target label is compatible with each retained order. No
individual-level source records are released.

Besides the empirical population $P$, we report
$P_{.01}=0.99P+0.01\U$, where $\U$ is uniform on the stated profile--label
space. This fixed smoothing has full support and gives a direct
example covered by Proposition~\ref{prop:local}. It is a diagnostic
regularization of the population, not an estimate that the data-generating
law has that contamination fraction.

\begin{table*}[t]
\centering
\caption{All eight-feature comparisons, including ranking-feasible profiles and the full-support population $P_{.01}$. Active counts coordinates on which the contrast depends. Twelve cube-domain deterministic updates are identically zero and omitted. Numerical certificate endpoints are released per task.}
\label{tab:mediumdetail}
\small
\setlength{\tabcolsep}{4pt}
\begin{tabular}{lrrrrrrr}
\toprule
Data / rules / domain & Tasks & Med. $b$ & Active & $\rho_{\rm wc}$ & $\rho_P$ & $\rho_{P_{.01}}$ & Max. $\rho_{P_{.01}}$ \\
\midrule
SUSHI / prob. / cube & 45 & 4 & 6 & 1.000 & 1.301 & 1.301 & 1.921 \\
SUSHI / prob. / ranking & 45 & 4 & 6 & 1.000 & 1.301 & 1.300 & 1.922 \\
SUSHI / det. / cube & 39 & 4 & 4 & 1.000 & 1.000 & 1.000 & 1.432 \\
SUSHI / det. / ranking & 39 & 4 & 4 & 1.000 & 1.000 & 1.000 & 1.432 \\
SAPA / prob. / cube & 25 & 4 & 6 & 1.000 & 1.260 & 1.262 & 2.435 \\
SAPA / det. / cube & 19 & 4 & 3 & 1.000 & 1.000 & 1.000 & 1.489 \\
\bottomrule
\end{tabular}
\end{table*}

\begin{table*}[t]
\centering
\caption{Primary-split larger-tree comparisons, conditional on 500 complete calibration records. Median variance ratios for unbiased scores; imputation uses MSE at $n=1000$. Oracle entries use upper certificates on the supplied menu, not all forms. Ordinal form scores use tree-threshold categories. Each task uses the fixed primary split; repeated-split diagnostics are reported separately.}
\label{tab:fulldetail}
\small
\setlength{\tabcolsep}{4pt}
\begin{tabular}{lrrrrrrr}
\toprule
Data & $b$ & Tasks & Prefix & Form scores & Oracle lower & Oracle upper & Imputation \\
\midrule
SUSHI & 3 & 45 & 2.613 & 2.607 & 2.532 & 2.532 & 0.948 \\
SUSHI & 6 & 45 & 3.765 & 3.403 & 3.184 & 3.184 & 0.835 \\
SAPA & 3 & 25 & 2.453 & 2.303 & 2.255 & 2.255 & 1.031 \\
SAPA & 6 & 25 & 2.758 & 2.438 & 2.286 & 2.286 & 1.030 \\
SAPA-ordinal & 3 & 25 & 2.393 & 2.232 & 2.162 & 2.162 & 1.388 \\
SAPA-ordinal & 6 & 25 & 3.088 & 2.778 & 2.389 & 2.389 & 1.235 \\
\bottomrule
\end{tabular}
\end{table*}


\subsection{EM Calibration from Restricted Pilot Records}\label{app:em}
For uniformly assigned forms, record only the form, stopping-prefix cell,
and label. Let $N_{\iota,c,y}$ be the resulting counts, and let $c_\iota(\ell)$ be
the observed prefix for leaf $\ell$. The joint leaf--label probabilities
$p_{\ell y}$ maximize the penalized observed-data log likelihood
\[
 \sum_{\iota,c,y}N_{\iota,c,y}\log\!\left(\sum_{\ell:c_\iota(\ell)=c}p_{\ell y}\right)
       +\tfrac12\sum_{\ell,y}\log p_{\ell y},
 \qquad p\ge0,\quad\sum_{\ell,y}p_{\ell y}=1.
\]
Form-assignment probabilities are known constants in this likelihood.
The half-count penalty prevents zero reference masses. The EM update
assigns each cell count among its compatible leaves in proportion to their
current probabilities, adds half a count to each leaf--label state, and
renormalizes. It uses unmatched records as coarse observations; it is not
allowed to use their hidden leaves. The implementation stops at maximum
coordinate change below $10^{-9}$, or 2,000 iterations, and reports the
achieved change and iteration count. A separate constrained-likelihood
optimizer checks the update on a small example.

The experiment freezes the 70 primary-split depth-six rules. Its population is the empirical law
of all records after the 800 training records, rather than the held-out
law of the conditional-calibration experiment. Each task has
500 iid simulated pilots, at each of $n_0=100,250,500$ respondents. Pilot
sizes share the beginning of the same random sequence to improve paired
comparisons. Every pilot respondent obeys the six-feature cap, and every
pilot record counts toward $N\in\{1000,2000,5000\}$.

We compare matched-path reference counts with EM reference estimates,
fitting the same prefix-score program afterwards. We also separate the
reference estimate from the score used to reuse the pilot. One pilot score
is completion-only weighting; a second is a globally unbiased prefix score
with uniform form probabilities, fixed in advance from the 800 historical
training records. The latter uses incomplete paths without fitting an
estimator to the same pilot that is being averaged.

For pilot mean $T_0$ and second-stage score mean $T_1$, use
$\lambda T_0+(1-\lambda)T_1$. We compare $\lambda=0$, $\lambda=n_0/N$, and
\[
 \lambda_{\rm prior}=
 \frac{\widehat V_1/(N-n_0)}{\widehat V_0/n_0+\widehat V_1/(N-n_0)},
\]
where both variance predictions come exclusively from the fixed historical
reference, before sampling the pilot. These weights are not selected with
evaluation data or with the realized pilot. Conditional on a pilot, the
exact risk is
\[
 \lambda^2(T_0-\Delta)^2+
       \frac{(1-\lambda)^2}{N-n_0}\Var_P(H_{\rm fitted}).
\]
The second-stage score remains unbiased for every fitted reference, so no
cross term remains. Averaging this expression over pilots yields the
reported risk estimates. Per-task Monte Carlo standard errors accompany
each comparison. These are risk estimates conditional on the frozen model
and empirical population, not independent replications of a deployment
study. All newly acquired records are capped at six features; with unit
feature and label costs the same comparison spends $7N$ total answers.

\begin{table*}[t]
\centering
\caption{Pilot-inclusive risk relative to branching, median across tasks. Each cell integrates the second stage exactly and averages 500 pilot replications per task. All new people obey the six-feature cap. Matched and EM estimate the reference; HT and stable specify the prespecified pilot score. Count, prior and discard specify reuse weights. The last column is the median per-task Monte Carlo standard error for EM/stable/prior, not uncertainty across populations.}
\label{tab:em}
\small
\setlength{\tabcolsep}{4pt}
\begin{tabular}{lrrrrrrrr}
\toprule
Data & $N$ & $n_0$ & Matched/HT & EM/HT & EM/stable & EM/prior & EM/discard & MC SE \\
\midrule
SUSHI & 1000 & 100 & 5.71 & 4.81 & 4.09 & 4.08 & 4.44 & 0.021 \\
SUSHI & 1000 & 250 & 7.32 & 6.04 & 4.30 & 4.21 & 5.37 & 0.054 \\
SUSHI & 1000 & 500 & 8.97 & 8.15 & 4.61 & 4.50 & 8.01 & 0.121 \\
SUSHI & 2000 & 100 & 5.33 & 4.37 & 4.05 & 4.02 & 4.21 & 0.012 \\
SUSHI & 2000 & 250 & 6.42 & 4.98 & 4.15 & 4.13 & 4.60 & 0.026 \\
SUSHI & 2000 & 500 & 6.97 & 6.06 & 4.29 & 4.25 & 5.34 & 0.054 \\
SUSHI & 5000 & 100 & 5.10 & 4.12 & 4.01 & 4.01 & 4.08 & 0.007 \\
SUSHI & 5000 & 250 & 5.93 & 4.38 & 4.07 & 4.06 & 4.24 & 0.012 \\
SUSHI & 5000 & 500 & 5.87 & 4.79 & 4.13 & 4.07 & 4.45 & 0.021 \\
SAPA & 1000 & 100 & 4.24 & 3.69 & 2.99 & 2.97 & 3.22 & 0.014 \\
SAPA & 1000 & 250 & 5.81 & 4.97 & 3.07 & 3.01 & 3.87 & 0.038 \\
SAPA & 1000 & 500 & 7.30 & 7.02 & 3.09 & 3.05 & 5.74 & 0.083 \\
SAPA & 2000 & 100 & 3.91 & 3.32 & 2.94 & 2.93 & 3.05 & 0.008 \\
SAPA & 2000 & 250 & 4.78 & 3.92 & 2.94 & 2.94 & 3.31 & 0.019 \\
SAPA & 2000 & 500 & 5.62 & 4.95 & 2.93 & 2.92 & 3.83 & 0.036 \\
SAPA & 5000 & 100 & 3.69 & 3.10 & 2.91 & 2.91 & 2.95 & 0.005 \\
SAPA & 5000 & 250 & 4.24 & 3.32 & 2.92 & 2.92 & 3.05 & 0.008 \\
SAPA & 5000 & 500 & 4.46 & 3.69 & 2.88 & 2.89 & 3.19 & 0.014 \\
\bottomrule
\end{tabular}
\end{table*}


At $N=1000,n_0=500$, changing only matched counts to EM while retaining
completion-weighted count-proportional reuse changes the median risk ratio
from 8.97 to 8.15 on SUSHI and from 7.30 to 7.02 on SAPA. Combining EM
with the stable historical pilot score and prior weight gives 4.50 and
3.05. The associated median per-task Monte Carlo standard errors are
0.121 and 0.083. At $n_0=100$ this last protocol gives 4.08 and 2.97, with
median errors 0.021 and 0.014. These paired comparisons identify the noisy
pilot score as a separate cost from reference estimation. Across all
105,000 fits, EM reached the specified coordinate tolerance in at most
629 iterations; none hit the 2,000-iteration cap. No claim of global
optimality over pilot sizes or allocation policies follows.

\subsection{Intervals for Frozen Unbiased Scores}\label{app:coverage}
With a fixed design fitted on separate historical or calibration records,
the audit observations $H_1,\ldots,H_n$ are iid conditional on that fitting
data and have mean $\Delta$. Use the ordinary sample variance
$s_H^2=(n-1)^{-1}\sum_i(H_i-\overline H)^2$ and the Wald interval
$\overline H\pm1.96s_H/\sqrt n$. For a finite menu with positive used
probabilities and positive variance, its conditional coverage tends to
$0.95$. This is an asymptotic statement; a small audit can miss rare large
importance scores. A finite-sample alternative uses bounded-score
concentration, at the cost of wider intervals.

We evaluate conditional coverage on the first five primary-split tasks
in each of SUSHI, binary SAPA and ordinal SAPA. The selection is by task
order, not observed coverage. Models and the complete 500-record
calibration are frozen. For each of $n=100,1000$, 2,000 iid empirical-population
audits compare branching, prefix scores, calibrated full-form scores,
and the fitted imputation comparator. Coverage is against the true gain of
that empirical population. Simulation standard errors are reported.
These results do not validate a Gaussian interval for an unknown imputation
bias: its interval generally targets the imputed mean instead of $\Delta$.
Nor do these single-stage intervals automatically cover an estimator that
reuses an adaptive pilot. With a prespecified pilot weight, the two-stage
variance must include both $\lambda^2s_0^2/r$ and
$(1-\lambda)^2s_1^2/(N-r)$; further conditions or a sequential interval are
needed to justify its normal approximation.

\begin{table*}[t]
\centering
\caption{Conditional coverage of nominal 95\% Wald intervals over the first five tasks per dataset, selected before inspecting coverage. Each task uses 2,000 iid audit replications with frozen training and calibration. Mean/min/max are descriptive summaries; the final column is the largest binomial Monte Carlo standard error within that row. Imputation intervals need not cover the target gain.}
\label{tab:coverage}
\small
\setlength{\tabcolsep}{4pt}
\begin{tabular}{lrrrrrr}
\toprule
Data & $n$ & Score & Mean & Min. & Max. & MC SE \\
\midrule
SUSHI & 100 & Branching & 0.944 & 0.939 & 0.949 & 0.005 \\
SUSHI & 1000 & Branching & 0.948 & 0.944 & 0.949 & 0.005 \\
SUSHI & 100 & Prefix & 0.947 & 0.942 & 0.954 & 0.005 \\
SUSHI & 1000 & Prefix & 0.948 & 0.945 & 0.955 & 0.005 \\
SUSHI & 100 & Form scores & 0.949 & 0.944 & 0.954 & 0.005 \\
SUSHI & 1000 & Form scores & 0.950 & 0.941 & 0.956 & 0.005 \\
SUSHI & 100 & Imputation & 0.934 & 0.879 & 0.955 & 0.007 \\
SUSHI & 1000 & Imputation & 0.797 & 0.242 & 0.951 & 0.010 \\
SAPA & 100 & Branching & 0.947 & 0.941 & 0.953 & 0.005 \\
SAPA & 1000 & Branching & 0.948 & 0.941 & 0.955 & 0.005 \\
SAPA & 100 & Prefix & 0.952 & 0.947 & 0.957 & 0.005 \\
SAPA & 1000 & Prefix & 0.949 & 0.945 & 0.953 & 0.005 \\
SAPA & 100 & Form scores & 0.942 & 0.940 & 0.948 & 0.005 \\
SAPA & 1000 & Form scores & 0.954 & 0.951 & 0.957 & 0.005 \\
SAPA & 100 & Imputation & 0.938 & 0.919 & 0.949 & 0.006 \\
SAPA & 1000 & Imputation & 0.782 & 0.618 & 0.946 & 0.011 \\
SAPA-ordinal & 100 & Branching & 0.944 & 0.939 & 0.948 & 0.005 \\
SAPA-ordinal & 1000 & Branching & 0.948 & 0.944 & 0.958 & 0.005 \\
SAPA-ordinal & 100 & Prefix & 0.945 & 0.942 & 0.950 & 0.005 \\
SAPA-ordinal & 1000 & Prefix & 0.951 & 0.949 & 0.958 & 0.005 \\
SAPA-ordinal & 100 & Form scores & 0.951 & 0.948 & 0.957 & 0.005 \\
SAPA-ordinal & 1000 & Form scores & 0.949 & 0.944 & 0.953 & 0.005 \\
SAPA-ordinal & 100 & Imputation & 0.265 & 0.018 & 0.503 & 0.011 \\
SAPA-ordinal & 1000 & Imputation & 0.000 & 0.000 & 0.001 & 0.001 \\
\bottomrule
\end{tabular}
\end{table*}


\subsection{Stratified Validation under Expected Cost}\label{app:validation}
To strengthen the relaxed-budget comparison, we also implement a two-phase
augmentation related to known-coarsening estimators and prediction-powered
inference \citep{robins1994,tao2020,ppi2024}. First observe a small fixed form
and the label, denoted $U$. Independently conditional on $U$, acquire the
remaining features with probability $\pi(U)>0$. For any frozen prediction
$m(U)$, the score
\[
          H=m(U)+\frac{R}{\pi(U)}\{Yd(X)-m(U)\}
\]
is conditionally unbiased given $(X,Y)$, and
\[
 \Var(H)=\Var(Yd)+\E\!\left[(\pi(U)^{-1}-1)(Yd-m(U))^2\right].
\]
For a fixed expected validation fraction, minimizing the second term uses
$\pi(u)$ proportional to the square root of the conditional squared
residual, clipped to the permitted probability range. This is the usual
Neyman allocation principle applied to augmentation residuals.

The same 500 complete calibration records fit smoothed cell means and
residual moments. Candidate initial forms are empty, the root alone, or
the root plus one other feature from a path form. Candidate mean validation
fractions are $0.05,0.1,0.2,0.4,0.7,1$, with $\pi\ge0.02$.
The initial form and fraction minimize estimated variance times estimated
cost using calibration alone. We then evaluate exact variance and expected
cost on the held-out population. Illustrative label costs are
$c_Y=1,10,100$ feature-answer units; these are sensitivity values, not
measured medical prices. Branching is charged for its actual mean path
length, not unused padding questions. The reported coefficient ratio is
\[
 \frac{[c_Y+|T|+(m-|T|)\E\pi(U)]\Var(H)}
      {[c_Y+\E\operatorname{depth}(X)]\Var(Yd)}.
\]
This is a large-budget, expected-cost comparison conditional on available
calibration. The validation decision depends on observed answers and can
request a complete profile, so it belongs to a relaxed two-phase workflow,
not to the hard-cap fixed-form class. It does not enforce an exact realized
total answer budget. It is reported separately from the finite
total-budget experiments to avoid conflating these constraints.

\begin{table*}[t]
\centering
\caption{Expected-cost coefficient ratios for stratified validation and complete-record averaging, relative to branching charged by actual mean path length. Calibration is historical and held fixed. Label cost $c_Y$ is in feature-answer units. Validation can request all features and depends on initial answers; it is outside the hard-cap fixed-form class.}
\label{tab:validation}
\small
\setlength{\tabcolsep}{4pt}
\begin{tabular}{lrrrrr}
\toprule
Data & $c_Y$ & Tasks & Stratified validation & Complete mean & Validation fraction \\
\midrule
SUSHI & 1 & 45 & 4.690 & 4.966 & 0.700 \\
SUSHI & 10 & 45 & 2.556 & 2.561 & 1.000 \\
SUSHI & 100 & 45 & 1.221 & 1.221 & 1.000 \\
SAPA & 1 & 25 & 3.167 & 3.572 & 0.688 \\
SAPA & 10 & 25 & 1.999 & 2.016 & 0.721 \\
SAPA & 100 & 25 & 1.144 & 1.144 & 1.000 \\
SAPA-ordinal & 1 & 25 & 3.683 & 3.815 & 0.397 \\
SAPA-ordinal & 10 & 25 & 2.081 & 2.068 & 0.691 \\
SAPA-ordinal & 100 & 25 & 1.148 & 1.148 & 1.000 \\
\bottomrule
\end{tabular}
\end{table*}

